% !TEX root = ../main.tex
\corrigesChapitre{6}{Décroissance radioactive}
\corrigesSouples

\begin{blocLivre}{cor-6-1}
\begin{corrige}{1}{Lire un extrait du diagramme (N,Z)}
\demarche{Les isotopes d'un élément sont sur une même ligne horizontale ; la couleur de chaque case indique le type de noyau.}
\begin{questions}
\item Le carbone a pour numéro atomique $Z=6$ : le noyau de carbone 14 contient 6 protons et $14-6=8$ neutrons :
\res{\ce{^{14}_{6}C}}.
\item Isotopes représentés : du carbone 10 au carbone 16 ; seuls \res{\text{le carbone 12 et le carbone 13}} sont stables.
\item Le carbone 11 est émetteur $\beta^+$ : il a trop de protons, il se situe à gauche des noyaux stables. Le carbone 14
est émetteur $\beta^-$ : il a trop de neutrons, il se situe à droite des noyaux stables.
\item $\res{\ce{^{14}_{6}C -> ^{14}_{7}N + ^{0}_{-1}e}}$ : $14=14+0$ et $6=7+(-1)$.
\end{questions}
\end{corrige}
\end{blocLivre}
\begin{blocLivre}{cor-6-2}
\begin{corrige}{2}{Écrire des équations de désintégration}
\demarche{Écrire la particule émise, puis déterminer $A$ et $Z$ du noyau fils par conservation et l'identifier par son numéro atomique.}
\begin{questions}
\item Au cours d'une réaction nucléaire, le nombre de nucléons $A$ et le nombre de charge $Z$ se conservent.
\item {\renewcommand{\arraystretch}{1.3}\begin{tabular}[t]{@{}l@{\hspace{3mm}}l@{\hspace{3mm}}l@{}}
A. & \ce{^{60}_{27}Co -> ^{60}_{28}Ni + ^{0}_{-1}e} & nickel 60\\
B. & \ce{^{18}_{9}F -> ^{18}_{8}O + ^{0}_{1}e} & oxygène 18\\
C. & \ce{^{241}_{95}Am -> ^{237}_{93}Np + ^{4}_{2}He} & neptunium 237\\
D. & \ce{^{3}_{1}H -> ^{3}_{2}He + ^{0}_{-1}e} & hélium 3
\end{tabular}}
\item La désintégration $\alpha$ de l'américium 241 : \res{\text{le noyau fils compte quatre nucléons de moins}}.
\end{questions}
\end{corrige}
\end{blocLivre}
\begin{blocLivre}{cor-6-3}
\begin{corrige}{3}{Identifier un type de radioactivité}
\demarche{Calculer les nombres manquants par conservation de $A$ et de $Z$, puis reconnaître la particule émise.}
\begin{questions}
\item {\renewcommand{\arraystretch}{1.3}\begin{tabular}[t]{@{}l@{\hspace{3mm}}l@{\hspace{8mm}}l@{\hspace{3mm}}l@{}}
A. & \ce{^{222}_{86}Rn -> ^{218}_{84}Po + ^{4}_{2}He} & B. & \ce{^{137}_{55}Cs -> ^{137}_{56}Ba + ^{0}_{-1}e}\\
C. & \ce{^{11}_{6}C -> ^{11}_{5}B + ^{0}_{1}e} & D. & \ce{^{232}_{90}Th -> ^{228}_{88}Ra + ^{4}_{2}He}\\
E. & \ce{^{99\text{m}}_{43}Tc -> ^{99}_{43}Tc + $\gamma$} & &
\end{tabular}}
\item A et D : \res{\alpha} ; B : \res{\beta^-} ; C : \res{\beta^+} ; E : \res{\gamma}.
\item Le photon $\gamma$ n'a ni charge ni nucléon : $Z$ ne change pas, l'élément reste le technétium. Le noyau
excité perd seulement de l'énergie.
\end{questions}
\end{corrige}
\end{blocLivre}
\begin{blocLivre}{cor-6-4}
\begin{corrige}{4}{Lire une courbe de décroissance}
\demarche{Lire la date où l'activité vaut la moitié de sa valeur initiale, puis compter les demi-vies.}
\begin{questions}
\item L'activité est le nombre de désintégrations par seconde ; elle s'exprime en \res{\text{becquerels (Bq)}}.
\item $A=\frac{400}{2}=\qty{200}{MBq}$ pour $\res{t_{1/2}\approx\qty{8}{jours}}$.
\item $\lambda=\dfrac{\ln2}{t_{1/2}}=\dfrac{\ln2}{\qty{8,0}{jours}}\approx\res{\qty{0,087}{jour^{-1}}}$.
\item On lit $\res{A\approx\qty{50}{MBq}}$. Par le calcul : $A=\frac{400}{2^3}=\qty{50}{MBq}$, car 24~jours correspondent
à trois demi-vies.
\end{questions}
\schema[corrige]{06-ex04}
\end{corrige}
\end{blocLivre}
\begin{blocLivre}{cor-6-5}
\begin{corrige}{5}{Compter les demi-vies}
\demarche{Exprimer chaque durée en nombre de demi-vies : l'activité est divisée par 2 à chaque demi-vie.}
\begin{questions}
\item 14~h : une demi-vie, \res{\qty{300}{MBq}} ; 20~h : deux demi-vies, \res{\qty{150}{MBq}} ; le lendemain à 8~h :
quatre demi-vies, $\frac{600}{2^4}=\res{\qty{37,5}{MBq}}$.
\item Deux jours correspondent à 8 demi-vies : il reste $\frac{1}{2^8}=\frac{1}{256}$, soit \res{\text{environ } 0{,}4~\%} de l'activité injectée.
\item L'activité devient rapidement négligeable : le patient est peu irradié après l'examen. En contrepartie, le
traceur doit être préparé peu de temps avant l'injection et utilisé rapidement.
\end{questions}
\end{corrige}
\end{blocLivre}
\begin{blocLivre}{cor-6-6}
\begin{corrige}{6}{Établir et utiliser la loi de décroissance}
\demarche{Dériver la fonction proposée, puis utiliser $N/N_0=\mathrm{e}^{-\lambda t}$.}
\begin{questions}
\item $\res{\ce{^{137}_{55}Cs -> ^{137}_{56}Ba + ^{0}_{-1}e}}$.
\item $\frac{\mathrm{d}N}{\mathrm{d}t}=N_0\times(-\lambda)\,\mathrm{e}^{-\lambda t}=-\lambda N$. À $t_{1/2}$, $\mathrm{e}^{-\lambda t_{1/2}}=\frac12$,
d'où $\lambda\,t_{1/2}=\ln2$ et $\lambda=\frac{\ln2}{30{,}0}=\res{\qty{2,31e-2}{an^{-1}}}$.
\item En 100~ans : $\frac{N}{N_0}=\mathrm{e}^{-\num{2.31e-2}\times100}\approx\res{9{,}9~\%}$.
\item $\frac{A_0}{A}=100$ pour $t=\dfrac{\ln100}{\lambda}=\dfrac{\ln100}{\num{2.31e-2}}\approx\res{\qty{2,0e2}{ans}}$.
\end{questions}
\end{corrige}
\end{blocLivre}
\begin{blocLivre}{cor-6-7}
\begin{corrige}{7}{Le détecteur de fumée}
\demarche{Convertir la demi-vie en secondes, puis utiliser $A=\lambda N$ et $m=\frac{N}{N_\mathrm{A}}\,M$.}
\begin{questions}
\item $\lambda=\dfrac{\ln2}{432{,}6\times\qty{3,156e7}{s}}\approx\res{\qty{5,08e-11}{s^{-1}}}$.
\item $N=\dfrac{A}{\lambda}=\dfrac{\num{3.7e4}}{\num{5.08e-11}}\approx\res{\num{7.3e14}}$ noyaux, soit
$m=\dfrac{N}{N_\mathrm{A}}\,M=\dfrac{\num{7.3e14}}{\num{6.02e23}}\times241\approx\res{\qty{2,9e-7}{g}}$, moins d'un microgramme.
\item $A=37\times\mathrm{e}^{-\frac{\ln2\times10}{432{,}6}}\approx\res{\qty{36,4}{kBq}}$ : l'activité ne diminue que de
1,6~\% en dix ans ; la source reste active pendant des siècles.
\item Les particules $\alpha$ sont arrêtées par le boîtier et par quelques centimètres d'air : pas d'exposition
externe. En revanche, une source détériorée pourrait contaminer l'environnement : l'inhalation ou l'ingestion
d'américium serait dangereuse, d'où une collecte spécifique.
\end{questions}
\end{corrige}
\end{blocLivre}
\begin{blocLivre}{cor-6-8}
\begin{corrige}{8}{Dater les restes d'un foyer préhistorique}
\demarche{Calculer le rapport $A/A_0$, le reporter sur la courbe, puis isoler $t$ dans la loi de décroissance.}
\begin{questions}
\item Un être vivant échange du carbone avec son milieu (respiration, photosynthèse, alimentation) : il renouvelle
le carbone 14 qui se désintègre, et sa proportion reste celle de l'atmosphère. Après la mort, les échanges cessent :
le carbone 14 n'est plus renouvelé et son activité décroît.
\item $\frac{A}{A_0}=\frac{0{,}060}{0{,}23}\approx0{,}26$ : on lit $\res{t\approx\num{11000}\text{ ans}}$.
\item $t=\dfrac{t_{1/2}}{\ln2}\ln\dfrac{A_0}{A}=\dfrac{\num{5700}}{\ln2}\ln\dfrac{0{,}23}{0{,}060}\approx\res{\num{1.1e4}\text{ ans}}$.
\item Au bout de \num{50000} ans, soit près de 9 demi-vies, l'activité est divisée par plus de 400 : elle devient
trop faible pour être distinguée de la radioactivité naturelle ambiante.
\end{questions}
\schema[corrige]{06-ex08}
\end{corrige}
\end{blocLivre}
\begin{blocLivre}{cor-6-9}
\begin{corrige}{9}{Le radon dans les habitations}
\demarche{Utiliser $N=A/\lambda$ avec $\lambda$ en \unit{s^{-1}}, puis $t=\ln f/\lambda$ pour une division par $f$.}
\begin{questions}
\item $\res{\ce{^{222}_{86}Rn -> ^{218}_{84}Po + ^{4}_{2}He}}$.
\item $\lambda=\dfrac{\ln2}{3{,}82\times86\,400}\approx\qty{2,10e-6}{s^{-1}}$, donc $N=\dfrac{300}{\num{2.10e-6}}\approx\res{\num{1.4e8}}$ noyaux.
\item $t=\dfrac{\ln10}{\ln2}\,t_{1/2}=3{,}32\times3{,}82\approx\res{\qty{13}{jours}}$.
\item Le radon est inhalé : les particules $\alpha$ émises par le radon et ses descendants irradient directement les
tissus des poumons (contamination interne). \res{\text{Aérer et ventiler}} les pièces limite sa concentration.
\end{questions}
\end{corrige}
\end{blocLivre}
\begin{blocLivre}{cor-6-10}
\begin{corrige}{10}{Simuler le hasard des désintégrations}
\demarche{Le programme répète, seconde après seconde, un tirage au sort pour chaque noyau restant.}
\begin{questions}
\item \code{random.random()} renvoie un nombre aléatoire entre 0 et 1 : le test est vrai avec la probabilité $p$, il
simule \res{\text{la désintégration d'un noyau pendant une seconde}}. La boucle se répète tant qu'il reste plus de la
moitié des noyaux initiaux.
\item Le nombre d'itérations, c'est-à-dire la durée, en secondes, au bout de laquelle la moitié des noyaux au moins
s'est désintégrée : c'est \res{\text{une estimation de la demi-vie}}.
\item Chaque désintégration est aléatoire : le nombre de noyaux désintégrés fluctue d'une exécution à l'autre. Avec
un plus grand nombre de noyaux, les fluctuations relatives sont plus faibles : la loi statistique est mieux suivie.
\item $N_k=N_0/2$ pour $(1-p)^k=\frac12$, soit $k=\dfrac{\ln2}{-\ln(1-p)}=\dfrac{\ln2}{-\ln0{,}95}\approx\res{13{,}5}$ secondes.
La simulation donne 14 secondes avec \num{50000} noyaux (premier entier supérieur à 13,5), et entre 12 et 15 secondes
avec 500 noyaux.
\end{questions}
\end{corrige}
\end{blocLivre}
\begin{blocLivre}{cor-6-11}
\begin{corrige}{11}{L'iode 131 et la thyroïde}
\demarche{Écrire la conservation de $A$ et de $Z$, puis utiliser $A(t)=A_0\,2^{-t/t_{1/2}}$.}
\begin{questions}
\item $\res{\ce{^{131}_{53}I -> ^{131}_{54}Xe + ^{0}_{-1}e}}$.
\item Au bout d'une demi-vie : \res{\qty{275}{MBq}} ; au bout de 30~jours :
$A=550\times2^{-30/8{,}02}\approx\res{\qty{41}{MBq}}$.
\item Les électrons sont arrêtés en quelques millimètres de tissu : ils irradient la thyroïde, où l'iode est
concentré, en épargnant le reste du corps. Un rayonnement $\gamma$, très pénétrant, traverserait la glande sans y
déposer l'essentiel de son énergie.
\item La thyroïde, saturée d'iode stable, ne fixe plus l'iode radioactif éventuellement inhalé ou ingéré, qui est
éliminé par l'organisme.
\end{questions}
\end{corrige}
\end{blocLivre}
\begin{blocLivre}{cor-6-12}
\begin{corrige}{12}{Préparer un examen de TEP}
\demarche{Remonter le temps avec $A_0=A\,\mathrm{e}^{\lambda\,\Delta t}$, puis appliquer $E=m\,c^2$ à la masse d'un électron.}
\begin{questions}
\item $\res{\ce{^{18}_{9}F -> ^{18}_{8}O + ^{0}_{1}e}}$.
\item Entre 8~h~30 et 11~h, il s'écoule 150~min : $A_\text{départ}=250\times\mathrm{e}^{\frac{\ln2\times150}{109{,}7}}\approx\res{\qty{645}{MBq}}$.
\item La masse $2m_\mathrm{e}$ est convertie en deux photons d'énergies égales ; chacun emporte $E=m_\mathrm{e}c^2$, soit
$\num{9.11e-31}\times(\num{3.00e8})^2=\qty{8,20e-14}{J}$ ou encore $\frac{\num{8.20e-14}}{\num{1.602e-19}}\approx\qty{5,1e5}{eV}$ :
\res{\text{environ 511~keV}}.
\item 24~h correspondent à $\frac{24\times60}{109{,}7}\approx13{,}1$ demi-vies : il reste $2^{-13{,}1}\approx\res{\num{1e-4}}$,
soit 0,01~\% de l'activité injectée. Quelques heures de précaution suffisent ; le lendemain, l'activité est négligeable.
\end{questions}
\end{corrige}
\end{blocLivre}
\begin{blocLivre}{cor-6-13}
\begin{corrige}{13}{La plus vieille roche terrestre}
\demarche{Écrire l'équation-bilan avec $x$ particules $\alpha$ et $y$ électrons, puis exprimer le plomb formé en fonction de l'uranium restant.}
\begin{questions}
\item \ce{^{238}_{92}U -> ^{206}_{82}Pb + $x$ ^{4}_{2}He + $y$ ^{0}_{-1}e}. Conservation de $A$ : $238=206+4x$, d'où
$\res{x=8}$ ; conservation de $Z$ : $92=82+2x-y$, d'où $\res{y=6}$.
\item Chaque noyau d'uranium désintégré donne un noyau de plomb : $N_\mathrm{Pb}=N_{\mathrm{U},0}-N_\mathrm{U}$. Comme
$N_\mathrm{U}=N_{\mathrm{U},0}\,\mathrm{e}^{-\lambda t}$, $N_{\mathrm{U},0}=N_\mathrm{U}\,\mathrm{e}^{\lambda t}$ et
$\res{N_\mathrm{Pb}=N_\mathrm{U}\left(\mathrm{e}^{\lambda t}-1\right)}$.
\item $t=\dfrac{1}{\lambda}\ln\left(1+\dfrac{N_\mathrm{Pb}}{N_\mathrm{U}}\right)=\dfrac{\num{4.468e9}}{\ln2}\ln(1{,}979)\approx\res{\num{4.40e9}\text{ ans}}$.
\item Le carbone 14 n'est pas présent dans un minéral ; et sa demi-vie (\num{5700} ans) est bien trop courte : au bout de
4,4 milliards d'années, il n'en resterait aucun noyau.
\end{questions}
\schema[corrige]{06-ex13-cor}
\end{corrige}
\end{blocLivre}
\begin{blocLivre}{cor-6-14}
\begin{corrige}{14}{Se protéger en médecine nucléaire}
\demarche{Le débit de dose varie en $1/d^2$ ; chaque couche de 0,30~mm de plomb le divise par deux.}
\begin{questions}
\item \ce{^{99\text{m}}_{43}Tc -> ^{99}_{43}Tc + $\gamma$}. Le rayonnement $\gamma$, pénétrant, sort du corps et est détecté
par la caméra ; la demi-vie de 6,0~h laisse le temps de l'examen tout en limitant l'irradiation du patient.
\item La distance est multipliée par 5 : le débit de dose est divisé par 25, $\frac{1{,}4}{25}\approx\res{\qty{0,056}{mSv.h^{-1}}}$.
\item \qty{2,0}{mm} correspondent à $\frac{2{,}0}{0{,}30}\approx6{,}7$ couches de demi-atténuation : facteur
$2^{6{,}7}\approx\res{\num{1.0e2}}$.
\item Sans protection, la dose de 1~mSv est atteinte au bout de $\frac{1}{1{,}4}\approx\qty{0,71}{h}$, soit
\res{\qty{43}{min}} ; avec le protège-seringue, il faut cent fois plus de temps : \res{\text{environ 72~h}}. L'écran, la
distance et des gestes rapides sont indispensables.
\item 24~h correspondent à 4 demi-vies : $\frac{740}{2^4}\approx\res{\qty{46}{MBq}}$.
\end{questions}
\end{corrige}
\end{blocLivre}
\begin{blocLivre}{cor-6-15}
\begin{corrige}{15}{Une pile nucléaire pour Mars}
\demarche{Calculer $N$ pour 1,0~g, puis $A=\lambda N$ et la puissance $P=A\times E$.}
\begin{questions}
\item $\res{\ce{^{238}_{94}Pu -> ^{234}_{92}U + ^{4}_{2}He}}$.
\item $N=\frac{1{,}0}{238}\times\num{6.02e23}=\num{2.53e21}$ ; $\lambda=\frac{\ln2}{87{,}74\times\num{3.156e7}}=\qty{2,50e-10}{s^{-1}}$ ;
$A=\lambda N\approx\res{\qty{6,3e11}{Bq}}$.
\item $P=A\times E=\num{6.3e11}\times5{,}59\times10^{6}\times\num{1.602e-19}\approx\res{\qty{0,57}{W}}$ pour 1,0~g, soit
$\res{\qty{2,0e3}{W}}$ pour \qty{3,5}{kg}.
\item Rendement : $\frac{110}{\num{2.0e3}}\approx\res{5{,}5~\%}$. L'essentiel de l'énergie est perdu sous forme de chaleur,
mais la source fonctionne sans entretien pendant des années, loin du Soleil et la nuit.
\item $2^{-14/87{,}74}\approx\res{0{,}90}$ : il reste environ 90~\% de la puissance initiale. Avec le polonium 210, la
puissance serait divisée par deux tous les 138~jours, et la mission ne pourrait pas durer.
\end{questions}
\end{corrige}
\end{blocLivre}
